\section{From logarithmic polynomials to integer powers}
\label{sec:transfer}

\subsection{Why a resultant is the right arithmetic object}

The Gaussian integers \(\ZZ[\ii]\) form a unique factorization
domain. Every nonzero Gaussian integer has complex modulus at least
one. If \(f\) is a primitive polynomial over this ring and divides a
polynomial over \(\QQ(\ii)\), Gauss's lemma makes that divisibility
valid over \(\ZZ[\ii]\) as well.

For polynomials \(f,U\), the resultant vanishes exactly when they
have a common root. If \(f\) has leading coefficient \(a\) and
roots \(y_1,\ldots,y_d\), then
\[
\Res(f,U)=a^{\deg U}\prod_{\ell=1}^{d}U(y_\ell).
\]
Using a \emph{minimal} polynomial, rather than the whole binomial
that first defines \(y\), guarantees the nonzero product needed below.

\begin{lemma}[Uniform algebraic-value bound]\label{lem:resultant}
Let \(k\ge2\), let \(p/q>0\) be reduced, and suppose
\(|p/q-\pi^k|<1\). Put
\[
\theta=(p/q)^{1/k}>0,\qquad y=\ii\theta/2.
\]
For every \(n\), some \(h\) has \(U_{h,n}(y)\ne0\).
For each \(\eta>0\) and all sufficiently large \(n\),
this \(h\) satisfies
\begin{equation}\label{eq:alglower}
|U_{h,n}(y)|\ge
(2^kq)^{-m}\exp(-(k-1)(G+\eta)n).
\end{equation}
The threshold is independent of \(p,q\).
\end{lemma}
\begin{proof}
Dividing the matrix by \(\Pi_n\) preserves its invertibility.
It cannot send the nonzero vector \((1,y,\ldots,y^m)^{\mathsf T}\)
to zero, so some row \(U\) is nonzero at \(y\).

Let \(f\in\ZZ[\ii][Y]\) be a primitive minimal polynomial of \(y\)
over \(K=\QQ(\ii)\), of degree \(d\le k\), and let \(a\ne0\)
be its leading coefficient. It divides
\[
(2^kq)Y^k-\ii^k p
\]
over \(\ZZ[\ii]\). Thus \(a\mid2^kq\), in particular
\(|a|\le2^kq\). Every root of \(f\) belongs to
\(\{y\zeta:\zeta^k=1\}\). The error assumption implies
\(|y|<2\), so all these roots lie in the disk used in
Lemma~\ref{lem:height}.

Irreducibility of \(f\) and \(U(y)\ne0\) give
\(\Res(f,U)\ne0\). Its Gaussian-integrality yields
\[
1\le |a|^{\deg U}|U(y)|
\prod_{\ell=2}^{d}|U(y_\ell)|.
\]
Use \(\deg U\le m\), \(d\le k\), and the uniform height bound.
Since \(G+\eta>0\), replacing \(d-1\) by \(k-1\) only weakens
the result and gives~\eqref{eq:alglower}.
\end{proof}

No assumption of irreducibility of \(Y^k-\ii^kp/(2^kq)\)
has been made. If \(p/q\) is a rational power, or a binomial splits
over \(K\), the true degree \(d\) is simply smaller. In particular,
the argument covers every denominator and not just a generic
subsequence of approximations.

\subsection{The quantitative comparison}

\begin{proof}[Proof of Theorem~\ref{thm:main}]
Write \(\Delta_k=D-(k-1)G>0\). Choose \(\eta>0\) small enough that
\(\Delta_k-k\eta>0\). All estimates can be made simultaneous:
for \(n\ge n_0\), every row and every \(|Y|\le2\) satisfy
\[
|U_{h,n}(\lambda)|\le e^{-(D-\eta)n},\quad
|U_{h,n}(Y)|,\ |T_{h,n}(Y)|\le e^{(G+\eta)n}.
\]
Here \(n\) was required to be a multiple of \(L\). Accordingly choose
\begin{equation}\label{eq:nchoice}
n=L\left\lceil
\frac{m\log(2^kq)+\log2}
 {L(\Delta_k-k\eta)}
\right\rceil.
\end{equation}
For sufficiently large \(q\), this is at least \(n_0\).
The extra rounding is at most \(L\), a fixed constant.

Select the row in Lemma~\ref{lem:resultant}. The ratio of the small
remainder to its nonzero value at \(y\) is at most
\[
(2^kq)^m e^{-(\Delta_k-k\eta)n}\le\tfrac12.
\]
The divided-difference identity consequently gives
\[
|\lambda-y|
\ge\tfrac12(2^kq)^{-m}e^{-k(G+\eta)n}.
\]
If the divided difference were zero, that same identity would
contradict the preceding factor \(1/2\); hence no division by zero
is hidden here.

Substitution of~\eqref{eq:nchoice} yields
\[
|\lambda-y|\ge c_{k,m,\eta}
q^{-m(D+G)/(\Delta_k-k\eta)}.
\]
Indeed
\(\Delta_k-k\eta+k(G+\eta)=D+G\), which explains the numerator
in the final exponent.

Finally,
\[
\left|\pi^k-\frac pq\right|
=|\pi-\theta|\sum_{a=0}^{k-1}\pi^{k-1-a}\theta^a
\ge2\pi^{k-1}|\lambda-y|.
\]
Fractions outside \(|p/q-\pi^k|<1\), or with \(p/q\le0\),
need no special treatment. Letting \(\eta\) tend to zero proves
the irrationality-exponent assertion.
\end{proof}

\subsection{Why a common divisor helps twice}

Before division, the rates are \(D_0,G_0\). Dividing all coefficients
and remainders by the same \(\Pi_n\) changes them to
\[
D=D_0+\gamma,\qquad G=G_0-\gamma.
\]
Thus the numerator \(D+G\) is unchanged, while the useful gap grows by
\(k\gamma\):
\[
D-(k-1)G=D_0-(k-1)G_0+k\gamma.
\]
This is why a divisor must be applied consistently to the remainder,
the coefficients, and the divided difference. Reducing only one
height in the final formula would give an incorrect improvement.
